\documentclass[10pt,a4paper,twocolumn]{article}

\usepackage{amsmath}
\usepackage{graphicx}
\usepackage{booktabs}
\usepackage{tabularx}
\usepackage{algorithm}
\usepackage{algpseudocode}
\usepackage{hyperref}

\title{Modeling a More Complex Classification System}
\author{Ali Hussain}
\date{}

\begin{document}

\maketitle


\section{Proposed Pipeline}

The proposed pipeline is designed to classify two-dimensional data when the
training data does not already have class labels. Most of the work is completed
offline on a computer. Only a small lookup step is performed on the embedded
device.

First, the training features are scaled using min--max scaling. DBSCAN is then
used to discover groups in the unlabelled training data. A Gaussian kernel
density estimate (KDE) is fitted separately to each discovered class. Each
density map is normalised to its own peak and sampled onto a fixed grid. The
grid values are then converted to \texttt{uint8} so that each cell requires
only one byte.

The embedded device does not run DBSCAN or KDE. Those steps are already
completed offline. The device stores the byte tables, the training scaling
bounds, the grid size, and the reject threshold.

For a new query point, the device scales the two features using the stored
training bounds. It then finds the correct grid cell and reads one byte from
each class table. The class with the largest stored value is selected. If the
largest value is below the reject threshold, the output is UNKNOWN instead of
forcing the point into a class.

\begin{algorithm}[!bt]
\caption{Table-lookup classification}
\label{alg:lookup}
\begin{algorithmic}[1]

\Require Raw query point $(x,y)$
\Require One byte table $Q_c$ for each class
\Require Training scaling bounds
\Require Grid size $G$
\Require Reject threshold $\tau$
\Ensure Predicted class or UNKNOWN

\State Scale $x$ and $y$ using the stored training bounds

\State $col \gets \lfloor \tilde{x}G \rfloor$
\State $row \gets \lfloor \tilde{y}G \rfloor$

\State Clip $row$ and $col$ to $[0,G-1]$

\For{each class $c$}
    \State $v_c \gets Q_c[row,col]/255$
\EndFor

\State $\hat{c} \gets \arg\max_c v_c$

\If{$v_{\hat{c}} < \tau$}
    \State \Return UNKNOWN
\EndIf

\State $p_c \gets v_c / \sum_k v_k$

\State \Return $\hat{c}$ and distribution $p$

\end{algorithmic}
\end{algorithm}

\section{Application}

One possible application for this pipeline is a small wearable motion
classifier that recognises different movement patterns from two sensor
features.

For this proposed device, I would reserve approximately 4~KB of flash for the
classification model. The final lookup tables from this experiment require
only 768 bytes, plus 16 bytes for the four scaling bounds, so the model fits
comfortably inside this budget.

Labelled training data may not be available because the device could collect
motion data during normal use without a person manually identifying every
sample. DBSCAN is useful in this situation because it can discover groups
directly from the unlabelled feature data.

The motion classes may also have non-round shapes in the feature space.
Different movements can create curved or hollow regions rather than simple
groups centred around one average point. A density-based method is therefore a
better match than a nearest-centre method such as k-means.

\section{Method and Setup}

The experiment used a synthetic dataset containing 1030 two-dimensional
points. The dataset was generated using the fixed random seed

\[
20250903.
\]

It contained two crescent-shaped groups, one ring-shaped group, and 90
uniformly scattered background points.

The dataset was divided into three parts:

\begin{itemize}
    \item 60\% training,
    \item 20\% validation,
    \item 20\% test.
\end{itemize}

The training set contained 618 points. It was used for DBSCAN clustering, KDE
fitting, and creation of the stored lookup tables.

The validation set was used to select the KDE bandwidth. The test set was kept
unused during training and tuning and was used only once at the end to evaluate
the frozen model.

Both features were min--max scaled using bounds calculated only from the
training set. The same bounds were used later to scale the validation, test,
and query points.

DBSCAN was applied to the scaled training data using

\[
\varepsilon = 0.06
\]

and

\[
\texttt{min\_samples} = 8.
\]

The baseline DBSCAN run discovered three clusters containing 209, 198, and
178 points. Another 33 training points were labelled as noise.

A Gaussian KDE was fitted separately to each discovered class. Noise points
were excluded from the KDE fitting because they were not assigned to a
discovered class.

The KDE bandwidth was selected using the validation set. The tested values
were

\[
0.03,\ 0.04,\ 0.05,\ 0.06,\ 0.08,\ 0.10,\ 0.12.
\]

The highest validation mean log-likelihood was produced by

\[
h = 0.05,
\]

so this value was used for the final model.

Each normalised KDE surface was sampled onto a

\[
16 \times 16
\]

grid and stored using \texttt{uint8}. The reject threshold was

\[
\tau = 0.12.
\]


\end{document}